\section{Projective modules over a local ring}
\label{section-projective-local-ring}

\noindent
In this section we prove a very cute result:
a projective module $M$ over a local ring is free
(Theorem \ref{theorem-projective-free-over-local-ring} below).
Note that with  the additional assumption that $M$ is finite, this result is
Lemma \ref{lemma-finite-flat-local}.
In general we have:

\begin{lemma}
\label{lemma-projective-free}
Let $R$ be a ring.  Then every projective $R$-module is free if and only if
every countably generated projective $R$-module is free.
\end{lemma}

\begin{proof}
Follows immediately from
Theorem \ref{theorem-projective-direct-sum}.
\end{proof}

\noindent
Here is a criterion for a countably generated module to be free.

\begin{lemma}
\label{lemma-freeness-criteria}
Let $M$ be a countably generated $R$-module with the following property:
if $M = N \oplus N'$ with $N'$ a finite free $R$-module, then
any element of $N$ is contained in a free direct summand of $N$.
Then $M$ is free.
\end{lemma}

\begin{proof}
Let $x_1, x_2, \ldots$ be a countable set of generators for $M$.
We inductively construct finite free direct summands
$F_1, F_2, \ldots$ of $M$ such that for all $n$ we have
that $F_1 \oplus \ldots \oplus F_n$ is a direct summand of $M$
which contains $x_1, \ldots, x_n$. Namely, given $F_1, \ldots, F_n$
with the desired properties, write
$$
M = F_1 \oplus \ldots \oplus F_n \oplus N
$$
and let $x \in N$ be the image of $x_{n + 1}$.
Then we can find a free direct summand $F_{n + 1} \subset N$
containing $x$ by the assumption in the statement of the lemma.
Of course we can replace $F_{n + 1}$ by a finite free direct
summand of $F_{n + 1}$ and the induction step is complete.
Then $M = \bigoplus_{i = 1}^{\infty} F_i$ is free.
\end{proof}

\begin{lemma}
\label{lemma-projective-freeness-criteria}
Let $P$ be a projective module over a local ring $R$.  Then any element of $P$
is contained in a free direct summand of $P$.
\end{lemma}

\begin{proof}
Since $P$ is projective it is a direct summand of some free $R$-module $F$, say
$F = P \oplus Q$.  Let $x \in P$ be the element that we wish to show is
contained in a free direct summand of $P$.  Let $B$ be a basis of $F$ such that
the number of basis elements needed in the expression of $x$ is minimal, say $x
= \sum_{i=1}^n a_i e_i$ for some $e_i \in B$ and $a_i \in R$.  Then no $a_j$
can be expressed as a linear combination of the other $a_i$; for if $a_j =
\sum_{i \neq  j} a_i b_i$ for some $b_i \in R$, then replacing $e_i$ by $e_i +
b_ie_j$ for $i \neq j$ and leaving unchanged the other elements of $B$, we get
a new basis for $F$ in terms of which $x$ has a shorter expression.

\medskip\noindent
Let $e_i = y_i + z_i, y_i \in P, z_i \in Q$ be the decomposition of $e_i$ into
its $P$- and $Q$-components.  Write $y_i = \sum_{j=1}^{n} b_{ij} e_j + t_i$,
where $t_i$ is a linear combination of elements in $B$ other than $e_1, \ldots,
e_n$.  To finish the proof it suffices to show that the matrix $(b_{ij})$ is
invertible.  For then the map $F \to F$ sending $e_i \mapsto y_i$ for
$i=1, \ldots, n$ and fixing $B \setminus \{e_1, \ldots, e_n\}$ is an
isomorphism,
so that $y_1, \ldots, y_n$ together with $B \setminus \{e_1, \ldots, e_n\}$
form a basis for $F$.  Then the submodule $N$ spanned by $y_1, \ldots, y_n$
is a free submodule of $P$; $N$ is a direct summand of $P$ since $N \subset P$
and both $N$ and $P$ are direct summands of $F$; and $x \in N$ since $x \in P$
implies $x = \sum_{i=1}^n a_i e_i = \sum_{i=1}^n a_i y_i$.

\medskip\noindent
Now we prove that $(b_{ij})$ is invertible. Plugging $y_i = \sum_{j=1}^{n}
b_{ij} e_j + t_i$ into $\sum_{i=1}^n a_i e_i = \sum_{i=1}^n a_i y_i$ and
equating the coefficients of $e_j$ gives $a_j = \sum_{i=1}^n a_i b_{ij}$.  But
as noted above, our choice of $B$ guarantees that no $a_j$ can be written as a
linear combination of the other $a_i$.  Thus $b_{ij}$ is a non-unit for $i \neq
j$, and $1-b_{ii}$ is a non-unit---so in particular $b_{ii}$ is a unit---for
all $i$.  But a matrix over a local ring having units along the diagonal and
non-units elsewhere is invertible, as its determinant is a unit.
\end{proof}

\begin{theorem}
\label{theorem-projective-free-over-local-ring}
\begin{slogan}
Projective modules over local rings are free.
\end{slogan}
If $P$ is a projective module over a local ring $R$, then $P$ is free.
\end{theorem}

\begin{proof}
Follows from Lemmas \ref{lemma-projective-free}, \ref{lemma-freeness-criteria},
and \ref{lemma-projective-freeness-criteria}.
\end{proof}



